\documentclass{article}

% ============================================================
% INFR11269 Machine Learning Systems — Coursework Report
% Academic year 2026/27
%
% This coursework uses the official ICLR 2027 LaTeX style.
%
% IMPORTANT FORMAT RULES
% - Do not edit iclr2027_conference.sty.
% - Do not change page geometry, text width/height, base font size,
%   or line spacing.
% - Do not enable \iclrfinalcopy.
% - Normal content packages are fine if they do not alter the report format.
%
% Main-body page limit: six pages.
% References and Appendix do NOT count toward the main-body page limit.
% Reviewers are not obliged to read Appendix at all. Please put your main content
% in the main body.
% 
% AI policy:
% Academic integrity is an underlying principle of research and academic practice.
% All submitted work is expected to be your own, including any material you have
% generated yourself with the assistance of AI tools such as ELM, ChatGPT,
% Bing Chat, Claude, Google Gemini, or any other automated assistance.
% For all such material you must make prominent explicit acknowledgement and
% explain clearly what you have done in the AI declaration section.
% This includes using AI to directly generate text, code, figures or similar
% or to complete other tasks. In all cases, you remain entirely responsible
% for the material you produce.
% ============================================================

\usepackage{iclr2027_conference,times}

\usepackage{hyperref}
\usepackage{url}
\usepackage{booktabs}
\usepackage{graphicx}
\usepackage{amsmath}
\usepackage{amssymb}

% The official ICLR archive includes math_commands.tex. It is optional:
\input{math_commands.tex}

\title{Custom CUDA Kernels for GLM-ASR and RAG System Investigation}

% Keep this line unchanged.
\author{Anonymous coursework submission}

% Keep this COMMENTED OUT.
% Review mode is part of the coursework format and is checked automatically.
% \iclrfinalcopy

\begin{document}

\maketitle

% Course-specific template fingerprint. Keep this line unchanged.
\lhead{Under review as a report for UoE 26/27 Sem1 INFR11269: Machine Learning Systems}

% Keep this COMMENTED OUT.
% \begin{abstract}
% \end{abstract}

% ============================================================
% REPORT CONTENT GUIDANCE

% The structure below is RECOMMENDED, not mandatory.
% You may rename, merge, split, or reorder sections if another structure
% better communicates your investigation.

% Regardless of section titles, the report should make it easy to find:
%   1. the problem, motivation, and relevant system context;
%   2. the supplied baseline and your substantive technical work;
%   3. experimental methodology, metrics, and controlled comparisons;
%   4. results covering task/system quality AND relevant system cost;
%   5. analysis of why the results occur, limitations, and conclusions.
% Well organized report would result in better scores.

% Depth is preferred over implementing many unrelated features.
% ============================================================

\section{Introduction}
\label{sec:introduction}

This report covers INFR11269 coursework on machine learning \emph{systems}:
implementing the GPU kernels behind GLM-ASR-Nano speech recognition, then
investigating retrieval-augmented generation (RAG) built on that stack.
The ASR model routes every heavy operation through four CUDA tasks; PyTorch
holds tensors and orchestrates the graph. Task~1 implements rotary position
embedding (RoPE)---an element-wise kernel whose difficulty is index arithmetic
over \texttt{[batch, seq\_len, heads, head\_dim]} layouts, not floating-point
novelty.

The baseline is the supplied skeleton with empty blanks and PyTorch fallbacks
(\texttt{./run.sh --cuda none}). My work fills the kernels and measures
correctness (\texttt{make test1--4}) plus end-to-end transcription accuracy.
This section focuses on Task~1; later tasks and the RAG extension follow the
same quality--cost evaluation pattern.

\section{System and Technical Approach}
\label{sec:approach}

\subsection{Supplied baseline}
The homework provides \texttt{csrc/task*.cu} with numbered blanks,
\texttt{launch\_*} grid configuration, CPU references in \texttt{tests/}, and
\texttt{glm\_asr\_cuda/ops.py} which selects CUDA or PyTorch per kernel group.

\subsection{Task 1: RoPE kernel design}
RoPE rotates pairs \((v[i], v[i+h])\) before attention, with $h=\mathrm{rotary\_dim}/2$:
$\angle = \mathrm{pos}\cdot \mathrm{inv\_freq}[i]$,
$v[i]' = a\cos\angle - b\sin\angle$,
$v[i+h]' = b\cos\angle + a\sin\angle$.
One launch covers both \texttt{q} and \texttt{k}: threads $0,\ldots,\text{q\_pairs}-1$
touch \texttt{q}, then \texttt{k}. Each thread maps a global pair index to
\texttt{row}, pair index \texttt{i}, and token position
$\mathrm{pos} = (\mathrm{row}/\mathrm{num\_heads}) \bmod \mathrm{seq\_len}$.
Elements beyond \texttt{rotary\_dim} are untouched (audio encoder uses
\texttt{rotary\_dim}$=32$ of \texttt{head\_dim}$=64$).

Implementation choices: \texttt{ROPE\_THREADS}$=256$ blocks sized by
\texttt{ceil(total\_pairs / 256)}; in-place updates; \texttt{cosf}/\texttt{sinf}
computed per pair (no lookup table). Blanks~1.1--1.4 implement indexing and the
rotation; \texttt{launch\_rope} is unchanged.

\section{Experimental Methodology}
\label{sec:methodology}

\textbf{Hardware:} NVIDIA RTX A6000 (sm\_86, 84 SMs; test harness reports
${\sim}768$~GB/s peak DRAM bandwidth).

\textbf{Software:} \texttt{mls-asr} conda env (Python~3.11, PyTorch~2.10+cu128,
transformers~5.0); CUDA via \texttt{/usr/local/cuda}; repository
\texttt{edin-mls-26-autumn/hw1-asr}.

\textbf{Quality metrics:} per-case \texttt{max\_err} against CPU reference
(\texttt{make test1}); transcription string match on \texttt{test\_audio.wav}
(\texttt{./run.sh --cuda task1}).

\textbf{Cost metrics:} kernel latency (ms) and effective bandwidth (GB/s) printed
by the test harness; share of peak bandwidth for the memory-bound encoder case.

\textbf{Controls:} same GPU for all sweeps; rebuild with \texttt{make clean}
when changing \texttt{NVCCFLAGS}; default \texttt{-O3 -std=c++17 -lineinfo}
except for the \texttt{--use\_fast\_math} experiment.

\section{Results and Analysis}
\label{sec:results}

\subsection{Task 1 correctness and end-to-end check}
All four \texttt{make test1} cases pass with \texttt{max\_err} $\leq 2\times 10^{-7}$
(default flags). \texttt{./run.sh --cuda task1} yields 100\% transcription accuracy
(\emph{Concord returned to its place amidst the tents.}) with \texttt{rope=cuda}
and remaining groups on PyTorch.

\subsection{Report questions (Task 1 guide)}
\textbf{Q1 (encoder bandwidth).}
The harness counts bytes moved as
$2 \cdot (\text{q\_rows} + \text{k\_rows}) \cdot \text{rotary\_dim} \cdot 4$
(read $+$ write over the rotated prefix).
For the audio-encoder case ($B{=}1$, $\mathrm{seq}{=}1500$, $H_q{=}H_k{=}20$,
$\text{rotary\_dim}{=}32$): $\text{q\_rows}{=}\text{k\_rows}{=}30{,}000$,
so $2 \cdot 60{,}000 \cdot 32 \cdot 4 = 15.36 \times 10^6$~bytes (${\sim}14.6$~MiB)
per call. At ${\sim}0.026$~ms this is ${\sim}601$~GB/s, about \textbf{78\%}
of the quoted 768~GB/s peak---consistent with a memory-bound kernel that also
executes trigonometry and indexing.

\textbf{Q2 (precomputed $\cos/\sin$ table).}
A table of shape $[\mathrm{seq\_len}, h]$ for $\cos$ and $\sin$ would add
$2 \cdot \mathrm{seq\_len} \cdot h \cdot 4$~bytes read per pair
(two floats) on top of existing \texttt{q}/\texttt{k} traffic, while removing
\text{FLOPs} for \texttt{cosf}/\texttt{sinf}. For the encoder case
($\mathrm{seq\_len}{=}1500$, $h{=}16$) that is $192$~KiB of table data but
\textbf{+8~bytes/pair} on the critical path. At ${\sim}78\%$ of peak already,
extra reads likely \emph{slow} the kernel unless angles were reused many times
without touching \texttt{q}/\texttt{k} (they are not). Tables help when
trigonometry dominates and memory headroom exists; here arithmetic is cheap
relative to streaming $q$ and $k$.

\textbf{Q3 (one launch for $q$ and $k$).}
The model invokes RoPE \textbf{396 times} per transcription. Fusing $q$ and $k$
halves kernel-launch count versus separate launches, avoids an extra global
synchronisation point, and keeps one grid configuration hot in the driver. Cost
is slightly more complex indexing (blank~1.2), which is negligible compared
with launch latency (${\sim}$microseconds per launch) accumulated hundreds of times.

\subsection{Try this experiments (Task 1)}
Table~\ref{tab:rope-threads} sweeps \texttt{ROPE\_THREADS} on the encoder case
(default \texttt{NVCCFLAGS}). Throughput is flat near 76--79\% of peak for
block sizes 64--512; \texttt{1024} drops slightly (fewer blocks, less latency hiding).
\texttt{32} is slowest (about 76\% of peak): many small blocks increase launch and
scheduling overhead relative to useful memory work.

\begin{table}[t]
\centering
\caption{RoPE encoder-case bandwidth vs.\ \texttt{ROPE\_THREADS} (RTX A6000).}
\label{tab:rope-threads}
\begin{tabular}{@{}rrl@{}}
\toprule
\texttt{ROPE\_THREADS} & Time (ms) & Effective bandwidth \\
\midrule
32  & 0.026 & 581 GB/s, 0.76 of peak \\
64  & 0.026 & 598 GB/s, 0.78 of peak \\
128 & 0.025 & 606 GB/s, 0.79 of peak \\
256 & 0.026 & 601 GB/s, 0.78 of peak \\
512 & 0.026 & 599 GB/s, 0.78 of peak \\
1024 & 0.026 & 591 GB/s, 0.77 of peak \\
\bottomrule
\end{tabular}
\end{table}

\textbf{Tiny cases (${\sim}0.002$~ms).}
Micro-benchmarks move only hundreds of floats; time is dominated by \textbf{fixed
launch latency} and synchronisation, not by bytes transferred, so reported GB/s
is near 0\% of peak even when correct.

\textbf{\texttt{--use\_fast\_math}.}
Rebuilding with \texttt{NVCCFLAGS} including \texttt{--use\_fast\_math} leaves tiny and
decoder cases passing but \textbf{fails the audio-encoder case}
(\texttt{max\_err}${\approx}1.9\times 10^{-4}$ vs.\ tolerance $10^{-5}$).
Encoder tokens reach $\mathrm{pos}\uparrow 1499$ with $\mathrm{inv\_freq}[0]{=}1$,
so $\angle$ is large; hardware-style \texttt{cosf}/\texttt{sinf} approximations
drift from the CPU reference. Decoder angles stay smaller ($\mathrm{seq}{=}64$),
so fast math still passes. I restored default flags after the experiment
(\texttt{make clean} then \texttt{make}).

\section{Discussion and Conclusion}
\label{sec:conclusion}

Task~1 confirms that RoPE is memory-bound on realistic encoder shapes and that
correct indexing across batched $q/k$ layouts matters as much as the rotation
math. Launch configuration is secondary once block size is reasonable ($\geq 64$
threads). Fast math is unsafe for large angles---a systems lesson for any kernel
that applies transcendental functions per element.

Remaining ASR tasks (reductions, tiled GEMM, attention) dominate end-to-end time;
the RAG track will add retrieval and generation cost. Future work: complete
Tasks~2--4 with the same test-plus-transcription methodology, compare
\texttt{--profile} against \texttt{--cuda none}, then evaluate RAG quality and
latency with optional CUDA ASR integration.

% References are excluded from the main-body page limit.
\bibliography{references}
\bibliographystyle{iclr2027_conference}

% Put your Appendix here if needed. Comment out if you don't need it.
\appendix
\section{Task 1 reproduction commands}
\label{app:task1-repro}

\begin{verbatim}
conda activate mls-asr
export PATH=/usr/local/cuda/bin:$PATH
cd edin-mls-26-autumn/hw1-asr
make test1
./run.sh --cuda task1
# ROPE_THREADS sweep: edit csrc/task1_rope.cu, then make test1
# Fast math: make clean && make test1 \
#   NVCCFLAGS="-O3 -std=c++17 -lineinfo -diag-suppress 177 --use_fast_math"
\end{verbatim}

\section{Placeholder: later tasks and RAG}
Add \texttt{make test2--4}, full \texttt{./run.sh --profile} output, and RAG
experiments here or in the main body when complete. Reviewers are not obliged to
read appendices; keep essential claims in Sections~\ref{sec:results}--\ref{sec:conclusion}.

% Keep this line unchanged.
\newpage

% Academic integrity is an underlying principle of research and academic practice.
% All submitted work is expected to be your own, including any material you have
% generated yourself with the assistance of AI tools such as ELM, ChatGPT,
% Bing Chat, Claude, Google Gemini, or any other automated assistance.
% For all such material you must make prominent explicit acknowledgement and
% explain clearly what you have done in this section.
% This includes using AI to directly generate text, code, figures or similar
% or to complete other tasks. In all cases, you remain entirely responsible
% for the material you produce.
\section*{AI declaration}
I used AI coding assistants (including Cursor/ELM) to help run GPU experiments,
derive bandwidth accounting, and draft LaTeX text for this report. I reviewed
all kernel code, re-ran \texttt{make test1} and \texttt{./run.sh --cuda task1}
on my machine, and edited the final wording. I remain responsible for the
submitted work.

% Please put the information of your group here without changing the format.
% This page will be deleted when the report is distributed for reviewing.
\section*{Group information}
Group ID: XXX

Information of group member A:
\begin{itemize}
    \item Name: jacksu
    \item UUN: s2510156
\end{itemize}

Information of group member B:
\begin{itemize}
    \item Name: TBD (partner)
    \item UUN: TBD (partner)
\end{itemize}

% Keep this COMMENTED OUT, unless there is more than one member and you feel like
% it is necessary to point out the contribution. Please notice that members from
% the same group technically receive the same score for the report, regardless of
% the contribution.
% \section*{Contribution}

\end{document}
